% Estudo dirigido de Biblioteconomia
% Modelo do ÉguaTex (https://www.eguatex.com.br), em domínio público pela licença CC0 1.0:
% use, modifique e distribua sem pedir permissão e sem citar a origem.
%
% Compile com pdfLaTeX (o ÉguaTex faz isso sozinho). Arquivo único: main.tex.
% Cada questão discursiva tem linhas para a resposta (\fillwithlines)
% e uma resposta esperada no ambiente solution.

% ---------- versão do professor ou do aluno ----------
% Com \printanswers as respostas esperadas aparecem em caixas, antes das linhas.
% Troque por \noprintanswers para entregar a folha do aluno só com as linhas.
\documentclass[12pt,a4paper]{exam}
\printanswers

% ---------- idioma, fontes e página ----------
\usepackage[T1]{fontenc}
\usepackage[utf8]{inputenc}
\usepackage[brazilian]{babel}
\usepackage{lmodern}
\usepackage{microtype}
\usepackage[a4paper,top=2cm,bottom=2.4cm,left=2.5cm,right=2.5cm]{geometry}
\usepackage{xcolor}
\usepackage{array}
\usepackage[hidelinks]{hyperref}

\definecolor{petroleo}{RGB}{16,78,94}

% ---------- instruções de preenchimento ----------
% Os textos em cinza explicam o que escrever em cada parte.
% Troque \instrucoestrue por \instrucoesfalse para esconder todos de uma vez.
\newif\ifinstrucoes
\instrucoestrue
\newcommand{\instrucao}[1]{\ifinstrucoes{\color{gray}\small\itshape #1\par}\fi}

% ---------- dados do estudo dirigido: preencha aqui ----------
\newcommand{\universidade}{Universidade Federal Exemplo}
\newcommand{\curso}{Curso de Biblioteconomia}
\newcommand{\disciplina}{Organização e Representação da Informação}
\newcommand{\tema}{Estudo dirigido 2}
\newcommand{\assunto}{Representação descritiva e temática}
\newcommand{\professor}{Nome do professor ou da professora}
\newcommand{\entrega}{Entrega na aula seguinte}

\hypersetup{pdftitle={\tema}}

% ---------- aparência das questões e das respostas ----------
\qformat{\colorbox{petroleo}{\textcolor{white}{\bfseries\,\thequestion\,}}\enspace\textcolor{petroleo}{\bfseries\thequestiontitle}\hfill}
\renewcommand{\solutiontitle}{\noindent\textbf{Resposta esperada:}\enspace}
\SolutionEmphasis{\small}
\setlength{\linefillheight}{0.32in}
\colorfillwithlines
\definecolor{FillWithLinesColor}{gray}{0.6}

% ---------- cabeçalho e rodapé ----------
\pagestyle{headandfoot}
\firstpageheader{}{}{}
\runningheader{\footnotesize\textcolor{petroleo}{\disciplina}}{}{\footnotesize\textcolor{petroleo}{\tema}}
\footer{}{\small\thepage}{}

\begin{document}

% ======================================================================
% Identificação
% ======================================================================
\noindent{\small\MakeUppercase{\universidade}\enspace\textbullet\enspace\curso}\par
\noindent{\color{petroleo}\rule{\textwidth}{1.5pt}}\par
\noindent{\LARGE\bfseries\color{petroleo}\tema\par}
\noindent{\large\itshape\assunto\par}
\smallskip
\noindent\disciplina\par
\noindent Professor(a): \professor\par
\medskip
\noindent Nome: \makebox[0.55\textwidth]{\hrulefill}\hfill \entrega\par

% ======================================================================
% Orientações
% ======================================================================
\bigskip
\instrucao{Em duas ou três frases: objetivo, material de apoio e se a resposta é individual.}
\noindent\begin{tabular}{@{}!{\color{petroleo}\vrule width 3pt}@{\hspace{6pt}}p{\dimexpr\textwidth-9pt\relax}@{}}
  \small\textbf{Orientações.} Leia o texto-base e o capítulo indicado na bibliografia da disciplina. Responda com suas palavras, em frases completas, usando o espaço das linhas. As respostas são individuais e podem ser escritas a lápis.
\end{tabular}

% ======================================================================
% Texto-base
% ======================================================================
\bigskip
\instrucao{Texto curto, próprio ou citado com a fonte, que apoie as questões.}
\noindent{\bfseries\color{petroleo}Texto-base}\par
\smallskip
\noindent A representação da informação em bibliotecas tem duas faces. A representação descritiva registra os dados que identificam um documento, como título, autoria, edição, local, editora, data e descrição física, e define os pontos de acesso por nome e por título. A representação temática descreve o assunto do documento, por meio de códigos de classificação, como a CDD e a CDU, e de termos de indexação, como cabeçalhos de assunto e descritores de um tesauro. As duas juntas permitem que o usuário encontre o documento tanto pelo que ele é quanto pelo que ele trata.

% ======================================================================
% Questões
% ======================================================================
\bigskip
\instrucao{Escreva de quatro a seis questões que peçam explicação, comparação ou aplicação. Ajuste a altura de \textbackslash fillwithlines à extensão esperada: 3 cm para cerca de cinco linhas.}

\begin{questions}

\titledquestion{Representação descritiva e temática}
Com base no texto-base, diferencie a representação descritiva da representação temática e cite dois elementos de cada uma.
\begin{solution}
A representação descritiva identifica o documento pelos seus dados bibliográficos e físicos e define pontos de acesso por nome e título; exemplos: título e editora. A representação temática indica o assunto do documento; exemplos: o número de classificação na CDD e os descritores de um tesauro.
\end{solution}
\fillwithlines{4.4cm}

\titledquestion{A segunda lei de Ranganathan}
Explique a lei ``A cada leitor o seu livro'' e dê um exemplo de serviço de biblioteca que a coloca em prática.
\begin{solution}
A lei afirma que todo usuário tem direito a encontrar o material que atende à sua necessidade, sem distinção. Exemplos: o serviço de referência que orienta a busca de cada pessoa, a política de desenvolvimento de coleções baseada no perfil dos usuários e a oferta de obras em braile ou em áudio.
\end{solution}
\fillwithlines{4.4cm}

\titledquestion{Referência segundo a NBR 6023:2018}
Escreva a referência do livro com os dados abaixo, conforme a NBR 6023:2018.
\begin{center}\small
  \begin{tabular}{@{}ll@{}}
    Autoria: Ana Maria Silva & Título: Organização da informação \\
    Edição: 2.\textordfeminine & Local e editora: São Paulo, Editora Exemplo \\
    Ano: 2020 & Páginas: 180 \\
  \end{tabular}\par
  {\footnotesize(dados fictícios, criados para o exercício)}
\end{center}
\begin{solution}
SILVA, Ana Maria. \textbf{Organização da informação}. 2. ed. São Paulo: Editora Exemplo, 2020. 180 p.
\end{solution}
\fillwithlines{1.6cm}

\titledquestion{ISBN e ISSN}
Qual é a diferença entre o ISBN e o ISSN? Diga que tipo de documento cada um identifica.
\begin{solution}
O ISBN identifica um livro, uma edição específica de uma publicação monográfica, e tem 13 dígitos desde 2007. O ISSN identifica uma publicação seriada, como um periódico ou uma revista, com 8 dígitos, e vale para todos os fascículos do título.
\end{solution}
\fillwithlines{2.4cm}

\titledquestion{Dado, informação e conhecimento}
Diferencie dado, informação e conhecimento, com um exemplo ligado a uma biblioteca.
\begin{solution}
Dado é o registro bruto, sem contexto, como o número 245 em um registro. Informação é o dado organizado e com significado, como saber que o campo 245 do MARC 21 traz o título. Conhecimento é a informação assimilada por uma pessoa e usada para agir, como o catalogador que sabe preencher o registro corretamente.
\end{solution}
\fillwithlines{4.4cm}

\end{questions}

\end{document}
